\chapter{Технологический раздел}

\section{Средства реализации}

В качестве языка программирования, используемого при написании лабораторной работы, был выбран C++~\cite{cpp-lang}.
	
\section{Реализация алгоритмов}

На листингах~\ref{lst:compress.cpp}~--~\ref{lst:decompress.cpp} представлены реализации сжатия и разжатия произвольных документов.
\includelistingpretty
	{compress.cpp}
	{c++}
	{Реализация сжатия документа}

\includelistingpretty
	{decompress.cpp}
	{c++}
	{Реализация разжатия документа}

\section{Демонстрация работы программы}

На рисунке~\ref{img:demo} представлена демонстрация работы программы.
\includeimage
	{demo}
	{f}
	{H}
	{.7\textwidth}
	{Демонстрация работы программы}

\section{Тестирование программы}

В таблице~\ref{tbl:functional_test} представлены функциональные тесты для разработанной программы.
\begin{table}[ht!]
	\begin{center}
		\captionsetup{justification=raggedright,singlelinecheck=off}
		\caption{\label{tbl:functional_test} Функциональные тесты}
		\begin{tabularx}{\textwidth}{|c|X|X|}
			\hline
			№ & Входные данные & Выходные данные \\
			\hline
			1 & пустой файл & пустой файл \\
			\hline
			2 & aaaaa & сжатые <<aaaaa>> \\
			\hline
			3 & сжатые <<aaaaa> & aaaaa \\
			\hline
			4 & abcde & сжатые <<abcde>> \\
			\hline
			5 & сжатые <<abcde>> & abcde \\
			\hline
			6 & зашифрованный файл file.txt & файл file.txt \\
			\hline
			7 & файл file.bmp & сжатый файл enc.bmp \\
			\hline
			8 & сжатый файл enc.bmp & файл dec.bmp \\
			\hline
			9 & файл file.zip & сжатый файл enc.zip \\
			\hline
			10 & сжатый файл enc.zip & файл dec.zip \\
			\hline
		\end{tabularx}
	\end{center}
\end{table}